\section{ReactiveML Context and Core Constructs}
ReactiveML is the closest language-level incarnation of the control model
described in Section~\ref{sec:boussinot-control}. It offers dedicated syntax,
compiler support, and a runtime tailored to instants and delayed absence
\cite{reactiveml2005}. We do not reproduce the language itself here; instead,
we use it as a point of reference for the library-based approach of \tempo{}.

The remainder of this section therefore focuses on the core control constructs
and their synchronous properties, as they appear in \tempo{}'s library API.
\tempo{} exposes a small set of constructs that mirror the common vocabulary of
reactive control languages. Signals are the primary communication medium within
an instant. A signal is either present or absent during an instant; presence is
established by emission, while absence can only be observed at the end of the
instant. This delayed observation is central to deterministic behavior.

\paragraph{Semantic contract (informal)}
\tempo{} guarantees determinism at instant boundaries: for a fixed input stream,
the set of signals present at each instant and the values delivered to awaiters
are deterministic. What is not specified is the internal evaluation order of
tasks within an instant; relying on that order can break determinism when using
shared mutable state. The semantics therefore define \emph{what} is observable
per instant, not the precise scheduling order inside it.
We use \emph{instant boundary} and \emph{end of the instant} interchangeably.

\paragraph{Non-guarantees}
\tempo{} does not guarantee deterministic outcomes when programs depend on the
intra-instant evaluation order or shared mutable state (e.g., reference cells).
Such behaviors are outside the semantic contract and must be avoided or
controlled by static analyses.

\paragraph{Minimal vocabulary}
\begin{itemize}
\item \textbf{Instant}: logical time step where reactions occur.
\item \textbf{Presence/absence}: a signal is present if emitted in the instant.
\item \textbf{Guard}: a condition that blocks a task until a signal is present.
\item \textbf{Weak preemption}: termination at the end of the instant.
\item \textbf{Determinism}: observables per instant are uniquely determined by inputs.
\end{itemize}

\paragraph{Determinism statement (informal)}
Let $I$ be a fixed input sequence per instant. \tempo{} guarantees that, for each
instant $t$, the multiset of event-signal emissions and the values returned by
\texttt{await} and \texttt{await\_immediate} are uniquely determined by $I$,
assuming programs avoid shared mutable state that depends on intra-instant
ordering and aggregate combine functions are pure and order-insensitive. No
guarantee is made about the relative order of tasks within an instant.
\paragraph{Observable vs non-observable (summary)}
\begin{itemize}
\item \textbf{Observable:} which signals are present at each instant boundary,
  and the values delivered to \texttt{await} and \texttt{await\_immediate}.
\item \textbf{Not observable:} the order in which tasks run within the instant,
  or the ordering of side effects on shared mutable state.
\end{itemize}
\paragraph{Aggregate signals and \texttt{await\_immediate}}
Aggregate signals are delivered only at instant end; thus their determinism is
defined by the combined value, not by the order of emissions within the instant.
\texttt{await\_immediate} is restricted to event signals, which allows safe
same-instant resumption without introducing ambiguity in the delivered value.
For deterministic aggregation, the combine function should be pure and
order-insensitive (associative/commutative), since the intra-instant order of
emissions is not part of the semantics.

\paragraph{Observable trace (informal)}
An execution can be summarized by a per-instant trace $au(t)$ that records:
\begin{itemize}
\item \textbf{Event signals:} the unique value emitted in instant $t$, or
  \emph{absent} if no emission occurs.
\item \textbf{Aggregate signals:} the combined value delivered at instant end
  (or the initial accumulator if no emission occurred).
\item \textbf{Outputs:} values emitted on the output signal at instant end.
\end{itemize}
The trace intentionally omits intra-instant scheduling order. Two runs are
observationally equivalent if they produce the same $au(t)$ for all $t$.

\paragraph{\texttt{await} vs \texttt{await\_immediate}}
\texttt{await} always resumes at the next instant, even if the signal is already
present. If the event signal is present, \texttt{await\_\allowbreak immediate}
may resume in the current instant. This can be useful for immediate reactions,
but it is restricted to event signals to preserve determinism.

\begin{figure}[h]
\centering
\begin{tikzpicture}[
  box/.style=tempo,
  draw,
  rounded corners,
  align=left,
  minimum width=70mm,
  minimum height=8mm
]
\node[box] {
Observable per instant:\\
1) set of present signals\\
2) values delivered to awaiters\\
3) outputs emitted to the environment
};
\end{tikzpicture}
\caption{Observables that define determinism at instant boundaries.}
\end{figure}

\paragraph{Determinism in action}
For the same input stream, the observables above do not change. For example,
if the host emits input \texttt{1} at instant $t$ and \texttt{2} at instant
$t{+}1$, then the program below deterministically emits \texttt{1} at $t{+}1$
and \texttt{2} at $t{+}2$.
\begin{lstlisting}[style=tempo]
let program input output =
  let v = await input in
  emit output v
\end{lstlisting}

\begin{table}[h]
\centering
\begin{tabularx}{\linewidth}{lY}
\hline
Instant & Observable events \\ \hline
$t$ & input emits \texttt{1}; no output yet. \\
$t{+}1$ & output emits \texttt{1}; input emits \texttt{2}. \\
$t{+}2$ & output emits \texttt{2}. \\ \hline
\end{tabularx}
\caption{Timeline for the simple echo program.}
\end{table}

\paragraph{Mini-specification (informal)}
The behavior of a program can be described per instant:
\begin{enumerate}
\item At instant start, all signals are absent.
\item Emissions during the instant mark signals present; event signals accept
  at most one emission.
\item \texttt{await} resumes at the next instant; \texttt{await\_immediate}
  may resume in the current instant for event signals already present.
\item Absence is observed only at instant end; guard rechecks occur at the next
  instant.
\end{enumerate}

\paragraph{Intuitive timeline}
Consider two instants with a single signal \texttt{s}. In instant $t$, a task
emits \texttt{s} and another task awaits it. The awaiting task can only resume
at the \emph{beginning} of instant $t{+}1$, because the presence of \texttt{s}
becomes observable only after the instant ends. If \texttt{s} is never emitted
in instant $t$, that absence is observable only at the end of $t$, and a
guarded task can only retry in instant $t{+}1$. Observable events are therefore
the set of signals present at instant boundaries, not intermediate scheduling
steps within the instant.

Emissions mark a signal as present and provide a value. \tempo{} supports two
signal kinds: event signals accept at most one emission per instant, while
aggregate signals combine multiple emissions through an associative folding
function. Suspension is expressed through \texttt{await}, which resumes at the
beginning of the next instant, and \texttt{await\_immediate}, which may resume
in the current instant if an event signal is already present.

Guards and preemption encode control flow. Guarded execution activates a
behavior only when a guard signal becomes present. Weak preemption terminates a
behavior at the end of the current instant after a preempting signal is
emitted, allowing other concurrent tasks to observe the emission within the
same instant. These constructs preserve determinism by maintaining a single
global notion of instants and deferring absence decisions to instant boundaries.
The distinction between strong and weak restart patterns is often encoded by
explicit pauses: inserting a \texttt{pause} between successive activations
ensures that a reaction cannot retrigger within the same instant, which matches
the weak restart discipline commonly used in synchronous control.

Guarded execution is \emph{blocking}, so a sequential nesting of \texttt{when\_}
does not implement a disjunction. A disjunction is obtained by running
multiple guarded branches in parallel, each able to unblock independently.
Similarly, nested \texttt{watch} constructs compose as disjunctions: a body is
preempted if any watched signal is emitted in the current instant. This is
consistent with the control interpretation where multiple guards or preemption
sources represent alternative enabling or disabling conditions.

\begin{lstlisting}[style=tempo]
(* Disjunction of guards: run guarded branches in parallel. *)
parallel [
  (fun () -> when_ g1 (fun () -> body ()));
  (fun () -> when_ g2 (fun () -> body ()));
];

(* Disjunction of preemption: stop when s1 or s2 is emitted. *)
watch s1 (fun () -> watch s2 (fun () -> body ()))
\end{lstlisting}
